\section{Related Work}\label{sec:related}
\subsection{Classical sentiment classification}
Pang, Lee and Vaithyanathan~\cite{pang2002} compared supervised classifiers on
movie-review polarity, establishing sentiment classification as a useful test
of document-level language understanding. Their work motivates a linear-model
baseline for this project. Wang and Manning~\cite{wang2012} subsequently showed
the strength of simple lexical baselines with bigrams. Short phrases can preserve
local sentiment interactions that individual word counts miss, motivating our
unigram-and-bigram TF-IDF representation. Bigrams in our implementation are formed
after classical preprocessing and therefore do not necessarily correspond to
adjacent words in the original review.

\subsection{Distributed representations}
Mikolov et al.~\cite{mikolov2013} introduced efficient continuous word-vector
architectures, including Skip-gram. Such representations capture distributional
similarity rather than treating each word as an unrelated dimension. Maas et
al.~\cite{maas2011} connected word representations to sentiment analysis and
introduced the Large Movie Review Dataset used here. Our Word2Vec baseline learns
embeddings from project training reviews only and averages them into document
vectors. This simple aggregation tests a dense alternative to TF-IDF, although
it removes word order and does not explicitly model negation scope.

\subsection{Contextual encoders and efficient adaptation}
BERT~\cite{devlin2019} learns bidirectional contextual representations from
unlabeled text. DistilBERT~\cite{sanh2019} uses knowledge distillation to obtain a
smaller encoder suitable for more constrained computation. We choose its
six-layer uncased checkpoint as the project's deep model. LoRA~\cite{hu2022}
freezes pretrained weights and trains low-rank updates, reducing the number of
parameters requiring optimization. In our implementation, the classification
head is also trained in full. Work on BERT fine-tuning~\cite{sun2019} motivates
considering how long documents are truncated; we compare head-only and head-plus-tail
inputs instead of assuming that the first tokens always contain the final verdict.

The project evaluates these approaches within one IMDB experiment. Published
benchmark scores are not treated as directly comparable because dataset
preparation, tuning budgets and evaluation protocols differ.
